\documentclass{osucourse}
\course{4530}

\begin{document}

\CatalogDescription
\PrerequisitesAndExclusions

\begin{FollowupCourses}
Math 3589, Stat 4202.
\end{FollowupCourses}

\begin{LearningOutcomes}
By the end of this course, students should be able to:

\begin{itemize}
\tightlist
\item
  Understand the basic concepts in probability and statistics.
\item
  Compute probabilities and statistics of discrete and continuous
  distributions.
\item
  Comprehend the probabilistic methods needed to analyze and critically
  evaluate statistical models and arguments.
\item
  Recognize the importance of statistical ideas.
\end{itemize}
\end{LearningOutcomes}

\begin{Text}
\emph{Probability}, by Pitman, published by Springer, ISBN:
9780387979748
\end{Text}

\begin{TopicsOutline}
\textbf{I. Discrete probability.}

\begin{enumerate}
\def\labelenumi{\arabic{enumi}.}
\item
  First principles: outcome spaces, basic counting techniques, and
  partitions.
\item
  Venn diagrams and the inclusion-exclusion principle.
\item
  Conditional probability and independence; decision trees and Bayes'
  Theorem.
\item
  Discrete random variables; mass and generating functions; joint
  distributions.
\item
  Binomial, hypergeometric, geometric, negative binomial, and Poisson
  variables; applications and relationships.
\item
  Statistics on discrete variables.
\end{enumerate}

\textbf{II. Continuous probability}

\begin{enumerate}
\def\labelenumi{\arabic{enumi}.}
\setcounter{enumi}{6}
\item
  First principles: density functions, calculation of probabilities and
  statistics.
\item
  Moments and moment-generating functions.
\item
  Common distributions and their applications; exponential, gamma,
  uniform, normal.
\item
  The central limit theorem and normal approximation to the binomial
  distribution.
\item
  Relationships between the exponential, gamma, and Poisson
  distributions.
\item
  Hazard rates and survival functions.
\item
  Cumulative distribution functions, percentiles, and change of
  variables.
\item
  Joint distribution of continuous variables; independence and marginal
  distributions; density of a function of two variables
\end{enumerate}

\textbf{III. Statistics Material (using supplementary materials)}

\begin{enumerate}
\def\labelenumi{\arabic{enumi}.}
\setcounter{enumi}{14}
\tightlist
\item
  Chi-square distribution
\item
  t distribution
\item
  F distribution
\end{enumerate}
\end{TopicsOutline}

\begin{SuggestedSchedule}
\week{1}{Outcome spaces, counting, partitions.}
\week{2}{Venn diagrams; inclusion-exclusion.}
\week{3}{Conditional probability, independence, Bayes' Theorem.}
\week{4}{Discrete random variables; mass and generating functions.}
\week{5}{Binomial, hypergeometric, geometric, negative binomial variables.}
\week{6}{Poisson variables; statistics on discrete variables. \emph{Midterm I.}}
\week{7}{Density functions; probabilities and statistics in the continuous case.}
\week{8}{Moments and moment-generating functions.}
\week{9}{Exponential, gamma, uniform, and normal distributions.}
\week{10}{Central limit theorem; normal approximation to the binomial.}
\week{11}{Exponential, gamma, and Poisson relationships; hazard rates and survival. \emph{Midterm II.}}
\week{12}{Cumulative distribution functions, percentiles, change of variables.}
\week{13}{Joint continuous distributions; independence and marginals.}
\week{14}{Chi-square, t, and F distributions.}
\finalsweek{Final examination.}
\end{SuggestedSchedule}

\end{document}
